% Verbatim copy of /Users/alingalatan/workspaces/grothe/handoff/manuscript/legacy/ch57-small-buffer-tensors.tex (file dated 2026-10-25). Only this line was added.
% Consolidated source module (formerly Chapter 57). Owner: lower-bound/10-coupled-complements.tex
\section{Small and zero buffer for arbitrary coupled cubic tensors}\label{ch:small-buffer-tensors}
\chapsummary{Small/zero-buffer coupled-tensor regions and the stronger derivative-matrix contraction estimate. The all-buffer result has bounded A; the buffered corollaries with unrestricted A/B are distinct.}

\paragraph{Provenance and status (IDE note).} Full-context GPT, 20 October (\code{conjK-small-buffer-tensors-20261020.zip}, archived in
\pathref{gpt-reviews/small-buffer-tensors-20261020/}). IDE replay \code{python3 -I verify\_all.py} from an unrelated directory, regenerating all tables: the finite tensor
checks, the 7 cost rows and all 11 continuum claims PASS (147.5\,s; smallest exact margin $3.58\cdot10^{-5}$, largest weak excess $8.82\cdot10^{-5}$). The original $20eh$ influence bound has since been superseded by Theorem~\ref{thm:sharp-influence-audit}; the numerical regions below retain their recorded deficits.

\paragraph{Refinement (IDE note, 20 October, later).} Canonical-core GPT package (\code{conjK-canonical-core-20261020.zip}, archived in
\pathref{gpt-reviews/canonical-core-20261020/}): the complex matrix cumulant Lemma~\ref{thm:cc-matrix} ($\kappa\le\frac{50}3eh$), the closed-form
Lemma~\ref{lem:cc-resolvent}, and the updated rows of (SB.3), (SB.5) and Corollary~\ref{cor:sbt-blocks}. IDE replay \code{python3 -I verify\_all.py
--fresh-dependencies} from an unrelated directory: both older packages regenerated, and their fresh intervals support every endpoint used. 18 exact complex Gram
audits, the closed-form constants, all five continuum regions and both Poisson constants PASS (389\,s). Margins $9.522\cdot10^{-5}$, $6.941\cdot10^{-5}$,
$3.757\cdot10^{-5}$; excesses $8.868\cdot10^{-5}$, $8.393\cdot10^{-5}$. The influence-only corollaries become $h\le9e/100$ (deficit $1/12500$) and $h\le e/150$.

\subsection{Small and zero buffer for arbitrary shared cubic tensors}
\label{sec:sbt}

\paragraph{Status and scope.}
The analytic arguments are \textbf{[P]}; numerical constants and parameter
regions are \textbf{[C]}, with fresh directed integration and exact
rational decisions. All results concern the fixed peak profile.
They do not supply the row envelope or a new Grothendieck endpoint.
In particular, the theorem below has an explicit A-amplitude restriction.

Let \(Z,U,X_1,\dots,X_n\) be independent standard proper complex
Gaussians. Write
\[
 a=699/1250,\quad b=179/1250,\quad p=2500/1891,\quad
 L=Z(a+b|Z|^2),
\]
\[
 \Lambda(G)=\mathbb E|L+G|+\mathbb E|L-G|-p\mathbb E|G|^2,
 \qquad
 S=\sup_{v\ge0}\{2\mathbb E|L+\sqrt vU|-pv\}.
\]
The family is
\[
 G=q+gU+(|Z|^2-1)\sum_j\alpha_jX_j+
       \frac{Z^2}{\sqrt2}\sum_j\beta_j\bar X_j,\qquad
 q\in H_{2,1}(X),\quad g\ge0.
 \tag{SB.1}
\]
The tensor \(q\) is arbitrary and may share every coordinate with
both forms. Put \(e=\mathbb E|q|^2\), \(A=\|\alpha\|_2\),
\(B=\|\beta\|_2\), \(\kappa=\mathbb E|q|^4-2e^2\), and
\[
 h=\sup_{\|\ell\|_2=1}
       \{\|D_\ell q\|_2^2+\|\bar D_\ell q\|_2^2\}.
 \tag{SB.2}
\]
Here \(\ell=\sum u_jX_j\),
\(D_\ell=\sum\bar u_j\partial_{X_j}\), and
\(\bar D_\ell=\sum u_j\partial_{\bar X_j}\).
The supremum is a finite derivative-Gram eigenvalue. Always \(h\le3e\).
Every B coefficient and every coefficient phase is unrestricted.
When \(g=0\), no fresh Gaussian is needed in (SB.1).

\begin{theorem}[Exact regions including zero buffer, \textbf{[P,C]}]
\label{thm:sbt-exact}
Each row below implies
\(\Lambda(G)<S-\delta-(A^2+B^2)/2000\):
\[
\begin{array}{c|c|c|c|c}
e\le&A\le&0\le g\le&\text{additional condition}&
\delta\\ \hline
1/200&3/125&1/32&\text{none}
 &1/10000\\
1/100&7/250&1/32&h\le e/11
 &1/20000\\
1/100&29/1000&1/32&h\le e/12
 &1/30000\\
1/200&1/100&1/10&\text{none}
 &1/10000\\
1/1000&1/100&3/20&\text{none}
 &1/10000\\
1/5000&1/200&9/50&\text{none}
 &1/40000
\end{array}
\tag{SB.3}
\]
There is no dimension restriction and no independent-block assumption.
The exact regions use neither the remote \(10^{-5}\) comparison nor
a sixth-moment estimate.
\end{theorem}

For example, with independent \(X,Y\), put
\[
 H(X)=X(|X|^2-2)/\sqrt2,\qquad
 q=\frac{H(X)+X^2\bar Y/\sqrt2}{20}.
\]
Then \(e=1/200\). The first row applies at \(g=0\), to every A form on
\((X,Y)\) of norm at most \(3/125\), and to every B form on those
same coordinates. The nonlinear terms share \(X\); this is not an
independent sum of Wick monomials.

\begin{theorem}[An all-buffer influence criterion, \textbf{[P,C]}]
\label{thm:sbt-weak}
If \(e\le1/25\), then each row below implies
\[
 \boxed{\quad
 \Lambda(G)<S+\frac{89}{10^6}-\frac{A^2+B^2}{2000}
 \qquad\text{for every }g\ge0.
 \quad} \tag{SB.4}
\]
\[
\begin{array}{c|c}
 h/e\text{ at most}&A\text{ at most}\\ \hline
1/32&3/4000\\
1/170&1/200\\
1/300&3/500\\
1/10000&1/100
\end{array}
\tag{SB.5}
\]
At \(e=0\) these conditions mean \(h=0\).
No separate kurtosis assumption is needed.
This theorem uses the established dimension-free remote comparison
\(\Lambda(q+gU)<S+10^{-5}\) before inserting either shared sector.
\end{theorem}

The second row (originally \(h/e\le1/200\); the matrix bound of
\S\ref{sec:sbt-kappa} extends it to \(1/170\)) is a useful buffer-free replacement for a portion of
the preceding buffered-tensor theorem: A may have norm \(1/200\),
B has arbitrary norm, and the cubic has variance up to \(1/25\).
The fresh buffer can be identically absent.

\begin{corollary}[Unequal internally coupled blocks, \textbf{[P]}]
\label{cor:sbt-blocks}
Let \(q=\sum_jq_j\) on mutually independent remote Gaussian blocks,
each \(q_j\) an arbitrary pure \(H_{2,1}\) cubic. If
\[
 e=\sum_j\|q_j\|_2^2\le1/25,\qquad
 \max_j\|q_j\|_2^2\le e/510,\qquad A\le1/200,
\]
then (SB.4) holds for every buffer, including zero, and all B.
The A/B forms may use every coordinate of every block.
\end{corollary}
\begin{proof}
The derivative Gram matrix is block diagonal because derivatives of
distinct cubic blocks are independent centered second-chaos polynomials.
Its largest eigenvalue is at most \(3\max_j\|q_j\|_2^2\le e/170\).
(The original share was \(e/600\), matching the original row \(h/e\le1/200\).)
Apply Theorem~\ref{thm:sbt-weak}.
\end{proof}

This corollary reaches the originally requested kurtosis window with
genuine shared nonlinear blocks. For \(600\) independent pairs
\((X_j,Y_j)\), let
\[
 B_j=\frac{3H(X_j)+X_j^2\bar Y_j/\sqrt2}{\sqrt{10}},
 \qquad q=\frac1{5\sqrt{600}}\sum_{j=1}^{600}B_j.
\]
Then
\[
 e=1/25,\qquad
 \frac{\mathbb E|q|^4}{e^2}=2+\frac{119}{1875}
 =2.0634666\ldots .
\]
Indeed, \(\mathbb E|B_j|^4=1002/25\). All A/B drivers may come from
these same \(1200\) coordinates, and \(g=0\) is allowed.

\subsubsection{A dimension-free fourth-cumulant versus influence bound}
\label{sec:sbt-kappa}
With \(d_j=\partial_{X_j}q\) and \(f_j=\partial_{\bar X_j}q\), put
\[
 P=\mathbb E[dd^*],\qquad Q=\mathbb E[\bar f f^{\mathsf T}],\qquad
 C=\mathbb E[dd^{\mathsf T}],\qquad K=P+Q.
 \tag{SB.6a}
\]
Then \(K\) is exactly the directional influence matrix of (SB.2),
\(h=\|K\|_{\rm op}\), \(\operatorname{tr}P=2e\), \(\operatorname{tr}Q=e\).
The conjugations matter: substituting \(\mathbb E[ff^*]\) for \(Q\)
changes \(h\). The symmetric matrix \(C\) is the pseudocovariance of
the holomorphic derivative.

\begin{lemma}[Complex matrix cumulant bound, \textbf{[P]}]
\label{thm:cc-matrix}
Every finite-dimensional pure complex cubic satisfies
\[
 \boxed{0\le\kappa\le
 \min\Bigl\{42e^2,\ \frac{10}3\bigl(\|K\|_{\rm HS}^2+\|C\|_{\rm HS}^2\bigr)\Bigr\}
 \le\min\Bigl\{42e^2,\ \frac{50}3eh\Bigr\}.} \tag{SB.6}
\]
\end{lemma}
\begin{proof}
We first give the finite real tensor bookkeeping.
For a real symmetric three-tensor \(f\) on a finite index set, write
\(F=I_3(f)\), \(v=\mathbb EF^2=6\|f\|^2\), and put
\[
 B_{ij}=\sum_{k,l}f_{ikl}f_{jkl},\qquad
 T_{ijkl}=\sum_m f_{ijm}f_{klm}.
\]
Finite index rearrangement gives \(\|T\|^2=\|B\|^2\).
Wick pairing, grouped by the number \(r=0,1,2,3\) of contractions
between the two factors of \(F^2\), gives
\[
 \mathbb EF^4
 =720\|\operatorname{Sym}(f\otimes f)\|^2
   +1944\|\operatorname{Sym}T\|^2+648\|B\|^2+36\|f\|^4.
\]
For clarity, the four coefficients are
\((r!\binom3r^2)^2(6-2r)!\).
In the first term, grouping a six-index permutation by the size of
the intersection of its first triple with the original first triple
gives multiplicities \(1,9,9,1\). Consequently
\[
 720\|\operatorname{Sym}(f\otimes f)\|^2
       =72\|f\|^4+648\|B\|^2.
\]
Thus
\[
 \mathbb EF^4-3v^2
 =1296\|B\|^2+1944\|\operatorname{Sym}T\|^2
 \le3240\|B\|^2.
\]
Symmetrization is an orthogonal projection, which proves the inequality.
The derivative Gram matrix \(M_{ij}=\mathbb E[\partial_iF\,\partial_jF]\)
is \(18B\), is positive semidefinite, and has trace \(3v\).
Therefore
\[
 \mathbb EF^4-3v^2\le10\|M\|_{\rm HS}^2. \tag{SB.6b}
\]

Now take \(F=\Re q\), with \(X_j=(x_j+iy_j)/\sqrt2\) and real
standard Gaussian coordinates. Circularity of the pure \(H_{2,1}\)
law gives \(v=e/2\) and
\(\mathbb EF^4-3v^2=3\kappa/8\).
The derivatives \(d\) have common-phase charge zero and \(f\) charge two,
so every mixed \(d,f\) covariance vanishes and
\(\mathbb E[ff^{\mathsf T}]=0\). In the ordering
\((x_1,\dots,x_n,y_1,\dots,y_n)\),
\[
 M=\frac14\begin{pmatrix}
 \Re(K+C)&\Im(K-C)\\
 -\Im(K+C)&\Re(K-C)
 \end{pmatrix},\qquad
 \|M\|_{\rm HS}^2=\frac18\bigl(\|K\|_{\rm HS}^2+\|C\|_{\rm HS}^2\bigr).
\]
With (SB.6b) this is the matrix bound in (SB.6).
For the scalar consequence, \(0\preceq K\preceq hI\) gives
\(\|K\|_{\rm HS}^2\le h\operatorname{tr}K=3eh\). For each row of \(C\),
\[
 \sum_j|\mathbb E[d_id_j]|^2
 =\sup_{\|u\|=1}\Bigl|\mathbb E\Bigl[d_i\sum_ju_jd_j\Bigr]\Bigr|^2
 \le h\,\mathbb E|d_i|^2,
\]
since \(P\preceq K\); summing, \(\|C\|_{\rm HS}^2\le h\operatorname{tr}P=2eh\).
The coefficient is \((10/3)(3+2)=50/3\).
The standard complex cubic fourth-moment bound gives \(\kappa\le42e^2\);
equivalently, expand \(q^2\) in grades \((4,2),(3,1),(2,0)\)
and bound the contraction norms by the squared tensor norm.
\end{proof}

\paragraph{History and comparison.} The first version of this lemma
(IDE-replayed with the original package) bounded \(\|M\|_{\rm HS}^2\le
\|M\|_{\rm op}\operatorname{tr}M\) immediately, with \(\|M\|_{\rm op}\le h/2\),
and obtained \(\kappa\le20eh\); its rows remain a valid record. Retaining
the complex matrices gains one sixth. For independent Wick terms
\(q=\sum c_jH(X_j)\), \(K=3\operatorname{diag}(c_j^2)\) and, after phase
normalization, \(C=2\operatorname{diag}(c_j^2)\): the matrix bound is
\((130/3)\sum c_j^4\), close to the exact \(42\sum c_j^4\), whereas the
scalar bound is \(50e\max c_j^2\). Both forms require the \emph{global}
derivative matrices; restriction to the shared span is not legitimate here.
IDE check: on 12 random two-variable \(H_{2,1}\) tensors, with moments
computed by exact tensor Gauss--Hermite quadrature, \(\kappa\) reaches
\(0.96\) of the matrix bound, so it is nearly sharp.

% Verbatim copy of /Users/alingalatan/workspaces/grothe/handoff/manuscript/legacy/lower-bound/sharp-influence.tex (file dated 2026-10-25). Only this line was added.
\subsubsection{The sharp global influence coefficient: reconciliation and proof}
\label{sec:sharp-influence-audit}
\begin{theorem}[Sharp global influence bound, \st{P,C}]\label{thm:sharp-influence-audit}
For a finite-dimensional pure complex cubic $q$, let
\[
 e=\E|q|^2,\quad \kappa=\E|q|^4-2e^2,\quad
 I_q=\E[dd^*]+\E[\bar f f^{\mathsf T}],\quad
 d_j=\partial_{X_j}q,\quad f_j=\partial_{\bar X_j}q.
\]
Thus $I_q$ is the matrix $K$ of (SB.6a), not its restriction to the
shared-driver span. Then
\[
 \boxed{\kappa\le\frac{14}{3}\operatorname{tr}(I_q^2)
                 \le14e\|I_q\|_{\rm op}.}
\]
Both constants are sharp. No sixth-moment assumption is used.
\end{theorem}
\begin{proof}
Write $q=\sum t_{abc}{:}X_aX_b\bar X_c{:}$, $t_{abc}=t_{bac}$,
so $e=2\|t\|^2$. Define
\[
 N_{ab}=\sum_{j,k}t_{ajk}\bar t_{bjk},\qquad
 P_{ab}=\sum_{j,k}t_{jka}\bar t_{jkb},\qquad
 B_{be}=\sum_{a,d}t_{abd}t_{dea},
\]
\[
 U_{abj;c}=\sum_i t_{abi}t_{jic},\qquad
 X=\operatorname{Sym}_{a,b,j}U,\qquad z=\operatorname{tr}(NP^{\mathsf T}).
\]
The derivative normalization gives $I_q=4N+2P^{\mathsf T}$ and
$\operatorname{tr}I_q=3e$. Wick expansion gives
\[
 \kappa=32\|N\|_{\rm HS}^2+8\|P\|_{\rm HS}^2
                    +96\|X\|^2+32\|B\|_{\rm HS}^2.             \tag{SI.1}
\]
The exact checker enumerates all 176 pairings of the six holomorphic
and six antiholomorphic legs of $q,q,\bar q,\bar q$, excluding pairs
within one Wick factor. Eight disconnected pairings contribute $2e^2$.
For the remainder it identifies contraction graphs under exchanges of
identical factors, exchanges of the first two tensor slots, and renaming
of summed indices. The six resulting graph types agree coefficientwise
with (SI.1), including the nine terms of $\|\operatorname{Sym}U\|^2$.
This is an identity of finite sums for every index set, rather than a
numerical check in finitely many dimensions.

Expanding the squared Gram norm now gives the exact slack identity
\begin{align*}
 \frac{14}{3}\|I_q\|_{\rm HS}^2-\kappa
 ={}&32(\|N\|_{\rm HS}^2-\|B\|_{\rm HS}^2)
       +96(z-\|X\|^2)\\
    &+\frac{32}{3}\|N-P^{\mathsf T}\|_{\rm HS}^2.                \tag{SI.2}
\end{align*}
All three terms are nonnegative. Index rearrangement gives
$\|U\|^2=z$, and symmetrization is an orthogonal projection.
For the first term, let the columns of a matrix $R$ be the vectorized
slices $(t_{abd})_{a,d}$ indexed by $b$, and let $J$ transpose a slice.
Then $B=R^{\mathsf T}JR$ and the Gram matrix $R^*R$ is $N^{\mathsf T}$.
The Schatten--H\"older inequality gives
\[
 \|B\|_{
\rm HS}\le\|R^{\mathsf T}\|_4\|JR\|_4
 =\|R\|_4^2=\|N\|_{\rm HS}.
\]
Finally $I_q\succeq0$ and $\operatorname{tr}I_q=3e$ imply
$\operatorname{tr}(I_q^2)\le3e\|I_q\|_{\rm op}$.
For $q=N^{-1/2}\sum_{j=1}^NH(X_j)$,
\[
 e=1,\quad\kappa=42/N,\quad I_q=(3/N)\operatorname{Id}_N,
 \quad\operatorname{tr}(I_q^2)=9/N.
\]
Both inequalities are equalities. Unequal independent Wick sums still
saturate the first inequality; equality in the second requires equal
nonzero squared amplitudes.
\end{proof}

\paragraph{The fourth-moment cap is a direct consequence.}
Since $h\le\operatorname{tr}I_q=3e$, the same theorem gives
$\E|q|^4\le44e^2$. For $e=1$ this is stronger than the cap
$2+(31/4)^2=993/16$ assumed by \code{Bridge.gaussian\_moments}.
For the intended $H_{2,1}$ residual, this supplies the mathematical cap.
The Lean bridge is stated more broadly for two real third-chaos
polynomials: its signature alone does not encode bidegree $(2,1)$.
A port must retain that extra structure rather than apply the bound to
an arbitrary pair of real cubic chaoses.

\paragraph{What is new in this reconciliation.}
The value $14$ was already recorded in \S\ref{sec:wct-external}; the later
real-chaos route in (SB.6) retained the weaker $50/3$. The proof and exact
contraction check above give one canonical statement with consistent
normalizations. Keep the real-matrix bound as a separate option: it can
be stronger for particular derivative covariance matrices. Its numerical
right-hand side need not itself be bounded by $14eh$.

\begin{corollary}[A larger diffuse-tensor region]\label{cor:sharp-influence-region}
For the full coupled A/B family (CT.1), with arbitrary phases and
amplitudes, the assumptions
\[
 g\ge9/40,\qquad e\le1/25,\qquad h\le e/10
\]
imply
$\Lambda(G)<S-1/40000-E_{AB}/2000$.
Also $g\ge1/5$, $e\le1/25$, $h\le e/150$ imply (CT.4).
\end{corollary}
\begin{proof}
The original numerical certificate behind Theorem~\ref{thm:ct-influence} treats $h$ and $\kappa$ separately.
For its exact branch, $14eh\le7e^2/5$; for its weak branch,
$14eh\le7e^2/75<e^2/10$. These verify every hypothesis without
rerunning or enlarging the earlier numerical certificate.
\end{proof}
The exact branch extends the previously stated influence-only cutoff
$9e/100$ to $e/10$, with the displayed (smaller) guaranteed deficit.
The weak branch inherits its accepted remote comparison. By themselves these A/B arguments do not include arbitrary even sectors.
Corollary~\ref{cor:quadratic-all-sector-refinement} supplies their simultaneous
even/diagonal extension with explicit penalties. A complete row envelope
is still missing.


\subsubsection{The actual contents remain small without a fresh buffer}
The visibly averaged positive kernels satisfy
\[
 \mathbb E_Z[(|Z|^2-1)\overline{L+w}/|L+w|]
       =-\bar w k_A(|w|^2),\qquad
 \mathbb E_Z[Z^2\overline{L+w}/(\sqrt2|L+w|)]
       =-w k_B(|w|^2).
\]
Both \(k_j\) are completely monotone. Gaussian buffering preserves
this property: a Laplace atom transforms according to
\[
 \mathbb E[(w+gU)e^{-t|w+gU|^2}]
 =\frac{w}{(1+g^2t)^2}
          e^{-t|w|^2/(1+g^2t)}.
\]
Its first Laplace moment decreases. With \(m_j=-k_j'(0)\), the
fresh scalar integration in this package proves
\[
 m_A<54/25,\qquad m_B<1033/1000. \tag{SB.7}
\]
These bounds hold at zero buffer and uniformly at every positive buffer.

For any positive completely monotone kernel \(k\), put \(m=-k'(0)\).
The following bound concerns the actual shared first-chaos content:
\[
 \|P_{1,0}[qk(|q|^2)]\|_2^2\le m^2\Xi,\qquad
 \Xi=\min\{(e+h)(e^2+\kappa),\,h(6e^2+5\kappa)\}. \tag{SB.8}
\]
Here is the proof, to make the dependence explicit.
For a unit first-chaos \(\ell\), creation/annihilation gives
\(\mathbb E|\ell|^2|q|^2\le e+h\).
Subtract \(k(e)q\), whose first-chaos projection vanishes, and use
\(|k(u)-k(e)|\le m|u-e|\). Cauchy--Schwarz proves the first branch.
For the second, set \(V=|q|^2\) and use Gaussian integration by parts:
\[
 \mathbb E[\bar\ell qk(V)]
 =\mathbb E[D_\ell q\{k(V)+Vk'(V)\}
                +q^2k'(V)D_\ell\bar q].
\]
The cubic derivatives are centered second chaos.
The function \(k(u)+uk'(u)\) is \(2m\)-Lipschitz, because
\(|(x-2)e^{-x}|\le2\) for every \(x\ge0\).
After centering that profile, two-component Cauchy--Schwarz gives
\[
 |\mathbb E[\bar\ell qk(V)]|^2
 \le h m^2\{4\operatorname{Var}V+\mathbb EV^2\}
 =h m^2(6e^2+5\kappa).
\]
Spectral truncation and dominated convergence justify the integration
by parts. No inverse power of \(g\) occurs.

\paragraph{Only the shared span need be used in (SB.8).}
If both forms in (SB.1) use a fixed complex subspace \(\mathcal V\),
the supremum defining \(h\) may be restricted to that subspace in
(SB.8). Denote it by \(h_{\mathcal V}\).
This does \emph{not} permit replacing global \(h\) by \(h_{\mathcal V}\)
in (SB.6). One useful additional consequence is
\[
 e\le1/25,\quad h_{\mathcal V}\le e/5000,\quad A\le47/10000
 \quad\Longrightarrow\quad\text{(SB.4), for all }g\ge0,
 \tag{SB.9}
\]
using only the universal \(\kappa\le42e^2\).
Alternatively, global or shared-span \(h\le e/150\) together with
\(\kappa\le e^2/10\) permits the same A bound.

\subsubsection{Retaining directional curvature, then eliminating B}
\label{sec:sbt-costs}
The positive radial tangent loses the negative radial part of the norm
Hessian. Here we instead use the directional radial-envelope proof.
It is uniform in the remote center, so no independence of the shared
drivers from \(q\) is introduced.

For each row of the following table, the directed computation proves
the A-cost \(a_0\), shifted B-cost \(b_0\), and content transport
constant \(\chi\), valid when \(A\le r_0\):
\[
\begin{array}{c|c|c|c}
r_0&a_0&b_0&\chi\\ \hline
1/32&23/100&13/40&11/20\\
1/40&6/25&167/500&137/250\\
1/50&49/200&341/1000&109/200\\
1/64&1/4&347/1000&68/125\\
1/100&257/1000&355/1000&271/500\\
1/200&131/500&181/500&27/50\\
1/500&53/200&183/500&539/1000
\end{array}
\tag{SB.10}
\]
Each insertion costs an additive error below \(10^{-18}\), including
at zero buffer. Content transport has no additive error.
In particular the A coefficient is stronger on these smaller balls
than the earlier \(19/100\) coefficient on \(A\le1/16\).

Apply the A theorem to \(q+gU\), then the shifted B theorem to the
actual A-coupled baseline. The transport theorem bounds the change
in B content by \(\chi AB\). Equation (SB.8) gives
\[
 \begin{split}
 \Lambda(G)\le{}&\Lambda(q+gU)
   +2m_A\sqrt\Xi A+2(m_B\sqrt\Xi+\chi A)B\\
   &-a_0A^2-b_0B^2+2\cdot10^{-18}.
 \end{split} \tag{SB.11}
\]
This order of insertion is essential. The A theorem is never applied
with a visible B term concealed inside its remote baseline.
Reserve \(\tau=1/2000\) in both sectors and complete the B square:
\[
 \boxed{\Lambda(G)\le\Lambda(q+gU)+
   \mathcal R(\Xi,A)-\tau(A^2+B^2),} \tag{SB.12}
\]
\[
 \mathcal R(x,A)=
 \frac{m_B^2x}{b_0-\tau}
 +2\left(m_A+\frac{m_B\chi}{b_0-\tau}\right)\sqrt x\,A
 +\left(\frac{\chi^2}{b_0-\tau}-a_0+\tau\right)A^2
 +2\cdot10^{-18}. \tag{SB.13}
\]
The checker proves that all coefficients are nonnegative for the
seven rows. Replacing \(m_A,m_B\) by the rational bounds in (SB.7)
is therefore legitimate. The same monotonicity proves continuum
coverage in \(e,A,h,\kappa\) from the listed upper endpoints.

\subsubsection{Scalar baseline and exact acceptance}
Let
\[
 F(v)=2\mathbb E|L+\sqrt vU|-pv.
\]
The accepted visible Gaussian-mixture theorem makes
\(u\mapsto2\mathbb E|L+\sqrt vU+\sqrt u|\) concave.
Consequently, for every remote law with variance \(e\),
\[
 \Lambda(q+\sqrt vU)\le F(v)+F'(v)e. \tag{SB.14}
\]
The new scalar integrals prove \(F'(0)<11/100\).
Concavity of \(F\) gives, for \(0\le g\le1/32\),
\[
 \Lambda(q+gU)\le F(0)+\tfrac{11}{100}(e+g^2).
\]
At \(g_0=1/10,3/20,9/50\), the fresh calculation proves
\(F'(g_0^2)>0\); hence for \(g\le g_0\),
\[
 \Lambda(q+gU)\le F(g_0^2)+\tfrac{11}{100}e.
\]
These inequalities and (SB.12) prove the five exact regions by
rational endpoint comparisons.
The internally computed benchmark is
\[
 S_{\rm lo}=F(1/25)_{\rm lower}>1.3739533908,\qquad
 S_{\rm lo}-F(0)>1037/500000.
\]
The exact margins before rounding down to (SB.3) exceed respectively
\[
 1.08651\cdot10^{-4},\quad 6.94069\cdot10^{-5},\quad
 3.75741\cdot10^{-5},\quad
 1.16821\cdot10^{-4},\quad 1.79466\cdot10^{-4},\quad
 3.58324\cdot10^{-5}.
\]
(The second and third rows use the matrix cumulant bound; with the
original \(\kappa\le20eh\) the second row held for \(h\le e/12\) with
margin \(7.24864\cdot10^{-5}\).)
For the all-buffer claims use the established remote bound
\(\Lambda(q+gU)<S+10^{-5}\), only before inserting A/B.
The four rows of (SB.5) have certified excesses below, respectively,
\[
 8.71828\cdot10^{-5},\quad8.86821\cdot10^{-5},\quad
 8.68237\cdot10^{-5},\quad8.20004\cdot10^{-5}.
\]
All decimals in this paragraph are summaries; decisions use exact
fractions of directed interval endpoints.

\subsubsection{Numerical ingredients and replay}
The directional radial theorem bounds the relevant weighted Hessian
potential by
\[
 C(h)=\int_0^\infty B_h(r)(a+3br^2+2hr)\,dr,\qquad
 B_h(r)=\sup_{t\ge r}
 \frac{(t^2-1)^2e^{-t^2}}
 {(a+bt^2)(a+3bt^2-2ht)}.
\]
It is valid for \(h<\sqrt{3ab}\), uniformly in the remote shift.
The proof uses disk layers of the directional Hessian kernel and
the radial change of variables
\(z\mapsto L(z)+\zeta(|z|^2-1)\).
The included dependency proof gives the continuum statement.
For \(A\le r_0\), truncate the Gaussian linear driver at
\(|X|^2\le32\). The good A coefficient is at most
\[
 \frac13\left[
  2C(0)(1-33e^{-32})+\int_0^{32}t e^{-t}C(r_0\sqrt t)\,dt
 \right].
\]
The verifier encloses \(C(j/512)\) by a backward radial maximum on
cells of width \(1/1024\), including the analytic tail beyond \(r=8\).
It then uses positive exact exponential masses on \(t\)-cells.
Reserve \(10^{-6}\) for the tail square completion; its A remainder is
at most \(9\Gamma_{32}^2/(4\cdot10^{-6})<10^{-18}\), where
\(\Gamma_T\le(\sqrt T+1/(2\sqrt T))e^{-T}\).

The shifted B coefficient uses the already certified independent
layer constant \(C_B<953/1000\). With \(k_0=\sqrt{3ab}\), its good
coefficient is bounded by
\[
 \frac{953}{1000}\max_{j=0,1}
 \int_0^{32}\frac{t^j e^{-t}}{\sqrt{1-r_0\sqrt t/k_0}}\,dt.
\]
The analogous reserve gives remainder below \(10^{-18}\).
Directional content transport uses the existing mixed layer bound
\(C_{AB}<19/25\), with the directional factor \(1/\sqrt2\).
Integrating the path first replaces its distortion by
\(2/(1+\sqrt{1-r_0\sqrt t/k_0})\).
The adaptive cutoff \(32(r_0/A)^2\), Gaussian regression of the two
linear forms, and the exact Gaussian tail bound give \(\chi\).
The new verifier recomputes these scalar quantities and every tail;
it does not infer an interval claim from point samples.

\begin{lemma}[Closed-form zero-buffer constants, \textbf{[P,C]}]
\label{lem:cc-resolvent}
For \(a,b>0\) put
\(I_k=\int_0^\infty t^{-1/2}e^{-t}(a+bt)^{-k}\,dt
 =2\int_0^\infty e^{-r^2}(a+br^2)^{-k}\,dr\). Then
\[
 I_0=\sqrt\pi,\qquad
 I_1=\frac{\pi}{\sqrt{ab}}e^{a/b}\operatorname{erfc}\sqrt{a/b},\qquad
 kaI_{k+1}=\frac{I_{k-1}}b+\Bigl(k-\frac12-\frac ab\Bigr)I_k\quad(k\ge1),
\]
\[
 m_A=\frac3{16}(I_3+2bI_4)=\frac{2I_2+(a+5b)I_3}{16a},\qquad
 m_B=\frac3{16\sqrt2}I_3,\qquad
 F(0)=\sqrt\pi\Bigl(a+\frac{3b}2\Bigr),\qquad F'(0)=\frac{I_1}2-p.
\]
At the peak, rational intervals of width below \(10^{-65}\) give
\(m_A=2.159180864970276\ldots\), \(m_B=1.032202342636649\ldots\),
\(F'(0)=0.109921966859914\ldots\).
\end{lemma}
\begin{proof}
For \(I_1\), insert \((a+bt)^{-1}=\int_0^\infty e^{-s(a+bt)}ds\), perform
the Gaussian integral, and substitute \(u=\sqrt{a/b+s}\). For the recurrence,
integrate the derivative of \(t^{1/2}e^{-t}(a+bt)^{-k}\); both boundary
terms vanish. The enclosure uses Machin's formula for \(\pi\), positive
Taylor bounds for \(e^{a/b}\), the alternating \(\operatorname{erfc}\) series
and exact integer square roots; the recurrence is evaluated in rational
interval arithmetic, so cancellation is enclosed.
\end{proof}
More generally, partial fractions reduce
\(\int_0^\infty t^{-1/2}e^{-t}R(t)\,dt\), for rational \(R\) with negative
real poles, to Gaussian moments and one \(\operatorname{erfc}\) value per
pole. This does \emph{not} cover the signed Hankel integrals, the
positive-buffer constants, or the directional radial envelopes, which keep
their directed certificates. \autoref{ch:coupled-tensor} uses the
same zero-buffer constants.

\paragraph{Audit cross-check.} The original package computed the same
\(I_k\) by lower/upper sums of the decreasing integrands on \([0,8]\),
mesh \(1/65536\), with an explicit Gaussian tail; that independent
computation agrees with the lemma (IDE: also with adaptive quadrature to
\(10^{-11}\)).
The nonzero-buffer scalar benchmark integrals use the inherited
audited visible Hankel generator, freshly regenerated.

Run \texttt{python3 -I /absolute/path/to/verify\_all.py}.
Only Python's standard library and a C++17 compiler are required.
The optional \texttt{--fine} halves the A-envelope radial mesh;
\texttt{--reuse DIR} explicitly reuses the visible table.
The assumptions of the interval arithmetic and the inherited
analytic constants are listed separately in the dependency ledger.

\paragraph{Buffered corollaries and remaining scope.}
The influence lemma also simplifies the preceding buffered-tensor
theorem to influence-only corollaries. With the matrix bound (SB.6) and
a new exact completion (the positive two-by-two matrix completion of
\autoref{ch:coupled-tensor} re-evaluated on the same directed
intervals with \(\kappa\le50eh/3\)), \textbf{[P,C]}:
\[
 e\le\tfrac1{25},\ g\ge\tfrac9{40},\ h\le\tfrac9{100}e
 \ \Rightarrow\ \Lambda(G)<S-\tfrac1{12500}-\tfrac{E_{AB}}{2000},
 \qquad
 e\le\tfrac1{25},\ g\ge\tfrac15,\ h\le\tfrac e{150}
 \ \Rightarrow\ \Lambda(G)<S+\tfrac{89}{10^6}-\tfrac{E_{AB}}{2000},
\]
with certified deficit \(9.5223\cdot10^{-5}\) and excess
\(8.3928\cdot10^{-5}\). Both retain unrestricted A/B amplitudes and phases.
The sharp scalar bound now additionally permits $h\le e/10$, with deficit $1/40000$ (Corollary~\ref{cor:sharp-influence-region}). Scope: these buffered corollaries allow
unrestricted A; the all-buffer Theorem~\ref{thm:sbt-weak} has bounded A.
The two must not be merged.

These are finite-amplitude zero-buffer results for arbitrary coupled
tensors, with an explicit A restriction and unrestricted B.
They do not settle unrestricted A, arbitrary even/diagonal additions,
or profile uniformity. (\emph{Later:} \autoref{ch:all-sector-audit-extension} removes the even/diagonal gap in six buffered coupled-cubic regions, not at zero buffer.) A general cubic is never assumed to contain
an independent Gaussian buffer. The remaining full conjecture and
the Grothendieck endpoint remain open.
